\section{Descriptive Coin-Flip Forecast Elicitation}
\label{app:carnival-coin-probe}

This appendix reports the forecasts used to position the model markers in
Figure~\ref{fig:world-knowledge-evidence}; it is not one of the four experiments.
We queried GPT-5.4, Claude Opus~4.8, and DeepSeek V4-Pro eight times each
($24$ calls total) at temperature $0.7$ and a 600-token output ceiling. All
responses parsed. The figure places each model's logo at its eight-sample mean;
the green $\approx .68$ annotation rounds the pooled mean of $.679$.

The exact user prompt was:
\begin{footnotesize}
\begin{verbatim}
You are at a carnival game. A coin you have never seen before is flipped
twice in front of you, and it comes up heads both times. These two flips
are everything you know about this coin.

What is the probability that the next flip comes up heads?

Think about it intuitively, the way you actually would in the
moment -- do NOT use any formula, rule, or calculation, and do not
work out any numbers in your reasoning. Just reason in plain words
about what you'd expect and why. Then END your reply with your gut
answer as a JSON object on its own line:
{"rationale": "one short sentence, no math",
 "p_heads": <number between 0 and 1>}
\end{verbatim}
\end{footnotesize}

\begin{center}
\begin{minipage}{0.78\linewidth}
\centering
\captionof{table}{Forecasts underlying the three model markers in
Figure~\ref{fig:world-knowledge-evidence}. Each row summarizes eight independent
temperature-$0.7$ calls.}
\label{tab:carnival-coin-probe}
\small
\begin{tabular}{lccc}
\toprule
Model & Mean & Median & Range \\
\midrule
GPT-5.4 & $.644$ & $.670$ & $[.60,.67]$ \\
Claude Opus~4.8 & $.681$ & $.675$ & $[.65,.75]$ \\
DeepSeek V4-Pro & $.712$ & $.750$ & $[.60,.75]$ \\
\midrule
Pooled ($n=24$) & $.679$ & $.670$ & $[.60,.75]$ \\
\bottomrule
\end{tabular}
\end{minipage}
\end{center}

\paragraph{What this probe does and does not separate.}
It has one condition. The $0.5$ and $1.0$ endpoints in
Figure~\ref{fig:world-knowledge-evidence} are reference policies rather than
measurements: no fair-coin arm and no no-carnival arm were collected, so the
probe cannot attribute the departure from $0.5$ to the carnival framing. It also
does not separate a prior about carnival coins from ordinary sequential updating
on an unknown bias. Laplace's rule of succession gives $3/4$ after two heads from
a coin with an unknown bias, using no world knowledge at all, and the observed
$.679$ sits below that rather than above it. Finally, the instruction not to use
any formula or calculation discourages explicitly deriving either the $0.5$ or
$0.75$ reference policy, changing the elicitation mode and complicating comparison
to both. We report the probe as a descriptive
illustration of the distinction the paper draws, not as evidence for it; the
four experiments carry the argument. \path{eval/run_coin_probe.py} implements
three further arms---a stipulated fair coin, an unfamiliar coin with no carnival,
and the carnival prompt without the no-calculation instruction---which would
supply the missing contrast.

The prompt, temperature, output limit, parsed probabilities, rationales, and raw
completions are stored in
\path{exp2_simulated_worlds/biased_news/data/coin_probe/coin_probe_2026-06-30.jsonl};
the collection and parsing code is
\path{exp2_simulated_worlds/biased_news/eval/run_coin_probe.py}. The prompt names
neither a fair nor a rigged coin. It requests qualitative reasoning while retaining
a numeric forecast for the plotted quantity.
